\documentclass[11pt,a4paper]{amsart}

\usepackage{enumitem}
\usepackage[margin=1in]{geometry}

\title{Unverified Findings Index for the RH Draft}
\author{Working ledger}
\date{April 2026}

\begin{document}

\maketitle

\section*{Purpose}

This file is the quarantine ledger for claims, repairs, and endgame steps that
are not yet ready for promotion into the canonical draft
\texttt{paper/proof\_of\_rh.tex}.

Each item should record:
\begin{enumerate}[nosep]
    \item its source in the paper,
    \item its current verification state,
    \item the missing dependency or objection, and
    \item the concrete condition for promotion.
\end{enumerate}

Status vocabulary:
\begin{center}
\texttt{open, heuristic, computational, blocked, rejected, promoted}
\end{center}

Promotion rule: only the coordinator should move material from this file into
\texttt{paper/proof\_of\_rh.tex}, and only after an independent adversarial
verification pass.

\section*{Queue}

\begin{description}[style=nextline,leftmargin=2.9cm,font=\normalfont\bfseries]

\item[UV-001]
\textbf{Title.} Calibrated size versus microscopic toy regime.

\textbf{Status.} open.

\textbf{Source in paper.} \texttt{rem:wip-calibration-small-u}, currently
starting at line 5500 of \texttt{proof\_of\_rh.tex}.

\textbf{Claim.} Replace the fixed comparison scale $x=x_0(D)$ by a variable
scale $x=x(m)$ so that the calibrated toy parameter $u$ is forced into the
microscopic regime needed by the toy-side direct proposition.

\textbf{Why unverified.} The draft does not yet prove the calibration chain,
pairing estimates, and denominator bounds uniformly in the shrinking-$x$
regime.

\textbf{Needed for promotion.} A uniform variable-$x$ version of the
calibration argument with explicit bounds showing $|u|<u_0(D)$ in the regime
used later.

\item[UV-002]
\textbf{Title.} Pair-like versus finite-core endgame.

\textbf{Status.} open.

\textbf{Source in paper.} \texttt{rem:wip-pairlike-finitecore}, currently
starting at line 5604 of \texttt{proof\_of\_rh.tex}.

\textbf{Claim.} The present toy anomaly is justified for a pair-like local core,
but the general finite-core branch should instead be driven by the first nonzero
odd jet of the corrected transverse scalar.

\textbf{Why unverified.} The manuscript has not yet closed the genuinely
finite-core branch or proved that the actual core is always pair-like in the
required perturbative sense.

\textbf{Needed for promotion.} Either a theorem reducing all relevant local
cores to the pair-like branch, or a full finite-core endgame stated and proved
at the level of the first nonzero odd jet.

\item[UV-003]
\textbf{Title.} Parity/projective factorization and the local good-patch
frontier.

\textbf{Status.} heuristic.

\textbf{Source in paper.}
\texttt{rem:wip-parity-projective-factorization-collision-blow-up}, currently
starting at line 12447 of \texttt{proof\_of\_rh.tex}.

\textbf{Claim.} The residual exact fixed-shear lane may factor through a
parity-sterile blown-up interaction layer, with the first surviving layer
controlled by a one-projective-parameter quotient scalar and a reduced
exceptional-locus analysis.

\textbf{Why unverified.} The edge-law package and the reduction on the
exceptional locus are presented as plausible targets rather than closed
theorems.

\textbf{Needed for promotion.} A proved edge-law statement for the actual
corrected defects plus a rigorous reduction or exclusion mechanism on the locus
$M=0$.

\item[UV-004]
\textbf{Title.} Explicit-pointwise bridge to the local good-patch detector.

\textbf{Status.} open.

\textbf{Source in paper.}
\texttt{rem:wip-explicit-pointwise-bridge-good-patch-detector}, currently
starting at line 21277 of \texttt{proof\_of\_rh.tex}.

\textbf{Claim.} Decide whether the current explicit highest-new pointwise source
$r=q^{(7)}$ is already detected by the local good-patch compensator.

\textbf{Why unverified.} The draft identifies the bridge question but does not
yet prove the equivalence or the failure of equivalence.

\textbf{Needed for promotion.} A theorem showing either that the good-patch
detector already sees the same-tower source or that any remaining obstruction is
necessarily nonlocal and provenance-sensitive.

\item[UV-005]
\textbf{Title.} Inverse-branch-gap endgame in the overlap regime on the strict
three-run lane.

\textbf{Status.} heuristic.

\textbf{Source in paper.}
\texttt{rem:wip-inverse-branch-gap-endgame-fixed-shear}, currently starting at
line 22788 of \texttt{proof\_of\_rh.tex}.

\textbf{Claim.} The overlap regime may be attacked through the inverse-branch
gap $\Delta(q)$, either by forcing a finite-order jet mismatch at an overlap
witness or by rigidifying to affine or translated transport.

\textbf{Why unverified.} The decisive exclusion theorem for the overlap regime
is not yet proved, and the endpoint propagation needed for the strongest rigid
alternative is still conditional.

\textbf{Needed for promotion.} A closed theorem excluding genuine overlap on the
strict three-run lane, with all hypotheses and endpoint propagation justified.

\item[UV-006]
\textbf{Title.} Same-parity compression of the residual affine lane.

\textbf{Status.} heuristic.

\textbf{Source in paper.}
\texttt{rem:wip-same-parity-compression-affine-lane}, currently starting at
line 26024 of \texttt{proof\_of\_rh.tex}.

\textbf{Claim.} The residual affine obstruction should reduce from arbitrary
nonadjacent recurrence to isolated same-parity sign descents in the intercept
sequence, with only specific middle-limited triples remaining dangerous.

\textbf{Why unverified.} The endpoint bookkeeping and final reduction to a fully
rigorous isolated-triple statement have not yet been completed.

\textbf{Needed for promotion.} A proved classification of the residual affine
lane in terms of explicit same-parity obstruction intervals and their exclusion.

\item[UV-007]
\textbf{Title.} Current endgame status.

\textbf{Status.} open.

\textbf{Source in paper.} \texttt{rem:wip-final-endgame-status}, currently
starting at line 26369 of \texttt{proof\_of\_rh.tex}.

\textbf{Claim.} The current contradiction closes only in the pair-like branch,
while a genuinely finite-core branch remains possible and would need a higher
odd-order obstruction.

\textbf{Why unverified.} This is a proof-state summary rather than a finished
theorem: the general finite-core endgame is still missing.

\textbf{Needed for promotion.} Closure of UV-002 or an explicit theorem showing
that the pair-like branch exhausts the relevant local configurations.

\item[UV-008]
\textbf{Title.} Asymptotic contradiction versus full RH.

\textbf{Status.} open.

\textbf{Source in paper.} \texttt{rem:wip-high-height-only}, currently
starting at line 26400 of \texttt{proof\_of\_rh.tex}.

\textbf{Claim.} Upgrade the asymptotic contradiction to a full proof by making
the high-height comparison effective, extracting an explicit threshold $T_*$,
and supplying a separate low-height bridge.

\textbf{Why unverified.} The needed constants, explicit threshold theorem, and
finite verification bridge are described but not yet established in effective
form.

\textbf{Needed for promotion.} Explicit toy-side and zeta-side constants,
effective conversion from $Q$ to height, a concrete $T_*$, and a rigorous bridge
covering all heights below that threshold.

\item[UV-009]
\textbf{Title.} Toy microscopic expansion.

\textbf{Status.} open.

\textbf{Source in paper.} \texttt{rem:wip-toy-microscopic-expansion},
currently starting at line 26692 of \texttt{proof\_of\_rh.tex}.

\textbf{Claim.} Make the toy-side quadratic expansion quantitative by recording
an explicit Taylor remainder depending only on the compact $d$-range $D$.

\textbf{Why unverified.} The proof is currently schematic and relies on
smoothness and symmetry without a written quantitative remainder estimate.

\textbf{Needed for promotion.} A uniform analytic expansion in $(u,s,d)$ with an
explicit remainder term and a clean dependence on $D$.

\end{description}

\end{document}
